\documentclass[acmtog, nonacm]{acmart}


\usepackage{booktabs} % For formal tables
\usepackage{xcolor}

\newcommand\EG[1]{{\color{ForestGreen!70!Yellow}{#1}}}

\hyphenation{Ben-che-kroun}
\newcommand\Revision[1]{{\color{black}{#1}}}
\usepackage{mathtools}

% TOG prefers author-name bib system with square brackets
\citestyle{acmauthoryear}
%\setcitestyle{nosort,square} % nosort to allow for manual chronological ordering
\newcommand\todoChangy[1]{{\color{purple}{#1}}}
\newcommand\todopyc[1]{{\color{blue}{Peter: #1}}}
\usepackage[dvipsnames]{xcolor}
\usepackage{amsmath}
% \usepackage{subfigure}
\usepackage{algpseudocode}
\usepackage{tabularx}
\usepackage{bm}
\usepackage{xcolor}
\usepackage{cancel}
\usepackage{calc}
\usepackage{accents}
\usepackage{tikz}
\usepackage{pifont}
\usetikzlibrary{backgrounds}
\usetikzlibrary{decorations.pathreplacing, arrows.meta, tikzmark}
\usetikzlibrary{calc}
\usetikzlibrary {shapes.geometric} 
\usepackage{wrapfig}
\usetikzlibrary {fit}

\usepackage{cleveref}
\usepackage{tikzscale}
\usepackage{graphicx}
% \usepackage{subcaption}
% \usepackage{graphicx}
\usepackage{units} % for nicefrac
\usepackage[percent]{overpic}
\usepackage{subfig}
\usepackage[export]{adjustbox}
\usepackage{rotating} % For rotated text

\usepackage{pgfplots}
\usepgfplotslibrary{groupplots}
\pgfplotsset{compat=1.18}


\usepackage[table,xcdraw]{xcolor}
\usepackage{booktabs} 
\usepackage{makecell}


\tikzset{every node/.style={font=\fontfamily{LinuxBiolinumT-TLF}\selectfont}}
\usepackage{bm}
\usetikzlibrary{spy}

\newcommand{\refsec}[1] {Section~\ref{#1}}
%\newcommand{\refeq}[1] {Eq.~(\ref{#1})}
\newcommand{\reffig}[1] {Figure~\ref{#1}}
\newcommand{\refalg}[1] {Algorithm~\ref{#1}}
\newcommand{\reftab}[1] {Table~\ref{#1}}
\newcommand{\refapx}[1] {Appendix~\ref{#1}}
\newcommand{\zhecheng}[1]{\textcolor{orange}{Zhecheng: #1}}
% \newcommand{\zhecheng}[1]{}

\newcommand{\mengfei}[1]{\textcolor{magenta}{Mengfei: #1}}
\newlength{\subcolumnwidth}
\newenvironment{subcolumns}[1][0.45\columnwidth]
 {\valign\bgroup\hsize=#1\setlength{\subcolumnwidth}{\hsize}\vfil##\vfil\cr}
 {\crcr\egroup}
\newcommand{\nextsubcolumn}[1][]{%
  \cr\noalign{\hfill}
  \if\relax\detokenize{#1}\relax\else\hsize=#1\setlength{\subcolumnwidth}{\hsize}\fi
}
\newcommand{\nextsubfigure}{\vfill}
% \usepackage{showframe}
\usepackage[ruled]{algorithm2e} % For algorithms
\renewcommand{\algorithmcfname}{ALGORITHM}
\SetAlFnt{\small}
\SetAlCapFnt{\small}
\SetAlCapNameFnt{\small}
\SetAlCapHSkip{0pt}

% Metadata Information
\acmJournal{TOG}
%\acmVolume{38}
%\acmNumber{4}
%\acmArticle{39}
%\acmYear{2019}
%\acmMonth{7}

% Copyright
%\setcopyright{acmcopyright}
%\setcopyright{acmlicensed}
%\setcopyright{rightsretained}
%\setcopyright{usgov}
%\setcopyright{usgovmixed}
%\setcopyright{cagov}
%\setcopyright{cagovmixed}

% DOI
%\acmDOI{0000001.0000001_2}

% Paper history
%\received{February 2007}
%\received{March 2009}
%\received[final version]{June 2009}
%\received[accepted]{July 2009}


% Document starts



\makeatletter
\newcommand{\showfontinfo}{
    Font: \f@family \f@series \f@shape, Size: \f@size
}
\makeatother
\begin{document}
% Title portion
\title{Supplementary Document for Lifting the Winding Number: \\ Precise Discontinuities in Neural Fields for Physics Simulation}

\author{Yue Chang}
\orcid{0000-0002-2587-827X}
\affiliation{%
  \institution{University of Toronto}
  \country{Canada}}
\email{changyue.chang@mail.utoronto.ca}


\author{Mengfei Liu}
\orcid{0000-0003-4989-6971}
\affiliation{%
  \institution{ University of Toronto}
  \country{Canada}}
\email{mengfei.liu@mail.utoronto.ca}


\author{Zhecheng Wang}
\orcid{0000-0003-4989-6971}
\affiliation{%
  \institution{ University of Toronto}
  \country{Canada}}
\email{zhecheng@cs.toronto.edu}

\author{Peter Yichen Chen}
\orcid{0000-0003-1919-5437}
\affiliation{%
  \institution{MIT CSAIL}
  \country{USA}}
\email{pyc@csail.mit.edu}


\author{Eitan Grinspun}
\orcid{0000-0003-4460-7747}
\affiliation{%
  \institution{ University of Toronto}
  \country{Canada}}
\email{eitan@cs.toronto.edu}

\maketitle

\renewcommand{\theequation}{A.\arabic{equation}} % Numbering format for the appendix
\setcounter{equation}{0} % Reset equation counter

\section{Integrating the generalized winding number}
Our implementation represents $\Gamma$ as a collection of polylines (allowing for more than one cut). Each polyline $ \mathbf{P} $ consists of $N$ segments. Each segment  $S_i$ is represented by a sequence of points $ \{\mathbf{p}_{i}\}_{i=1}^{N+1} $. 

Since we are going to calculate the integral of the generalized winding number from a partial curve, the segment set that is activated for winding number calculation is changing as the cutting extent $\alpha$ changes. We consider the fractional length of the curve being cut parameterized by $\alpha \in [0, 1]$. Let the segment length of the polyline $l_i$ be defined as $l_i = ||\mathbf{p}_{i+1} - \mathbf{p}_i||$. Then the total length of the polyline is $L = \sum_{i=1}^N l_i$. If $\alpha = 0.25$, this means including segments until the cumulative length reaches 25\% of the total curve length. 

The integration bounds are calculated by walking along the polyline until the total length reaches $\alpha L$. Practically, this means constructing $\Gamma^\alpha$ by sequentially including $\mathbf{p}_i$ from $\mathbf{p}_1$ to $\mathbf{p}_k$ for some $k$, stopping when the total length either meets or would exceed the desired fractional length $\alpha L$  by adding an extra point. If $\sum_{i=1}^k l_i = \alpha L$, then we have obtained our desired polyline that represents $\Gamma_\alpha$, otherwise we would need to partially include the next segment. This is achieved through a linear interpolation. Let $t$ represent the fraction of the last partial segment to be included, calculated as
\begin{align}
    t = \frac{\alpha L - \sum_{i=1}^k l_i}{L_k}.
\end{align}

Therefore, the new endpoint can be calculated as 
\begin{align}
\mathbf{p}^{'}_{k+1}=(1 - t)\mathbf{p}^{k} + t\mathbf{p}^{k+1}.
\end{align}

The new set of $k$ segments with points $ \{\mathbf{p}_{0} \dots \mathbf{p}_{k}, \mathbf{p}'_{k+1}\}_{\alpha} $ can then be used to calculate the generalized winding number $H(\mathbf{x}, \alpha)$ at a query point $\mathbf{x}$.

\begin{align}
\label{eq:discrete_gwn}
    H(\mathbf{x}, \alpha) =  \sum_{i=1}^{k-1} \frac{(\mathbf{x} - \mathbf{p}_{i})^\perp \cdot (\mathbf{p}_{i+1} - \mathbf{p}_{i})}{\|\mathbf{x} - \mathbf{p}_{i}\| \|\mathbf{x} - \mathbf{p}_{i+1}\|} +  \frac{(\mathbf{x} - \mathbf{p}_{k})^\perp \cdot (\mathbf{p}'_{k+1} - \mathbf{p}_{k})}{\|\mathbf{x} - \mathbf{p}_{k}\| \|\mathbf{x} - \mathbf{p}'_{k+1}\|}
\end{align}



\section{Handling Singularities}

The kernel function $K(\mathbf{x}, \mathbf{p}, r)$ is a spatially changing function centered on the end point of the curve $\mathbf{p}=\Gamma(0)$ or $\mathbf{p}=\Gamma(t)$, the radius $r$ of the kernel function is calculated by the distance between the end point and the length of the margin.
\begin{align}
    r &= \|\Gamma(0) - \Gamma(m)\| \quad \text{or} \quad r = \|\Gamma(\alpha) - \Gamma(\alpha-m)\|
\end{align}
We use the following formulation for the kernel function in our implementation:
\begin{equation}
K(\mathbf{x}, \mathbf{p}, r) =
\begin{cases}
\Revision{
   \frac{6\| \mathbf{x} - \mathbf{p} \|^2}{ r^2} \left( 1 - \frac{\| \mathbf{x} - \mathbf{p} \|}{r} \right)}, &  \text{for }  \Revision{0 \leq \| \mathbf{x} - \mathbf{p} \|  \leq \frac{r}{2} } \\[8pt]
 \Revision{1-\frac{1}{4}(2 - \frac{2\| \mathbf{x} - \mathbf{p} \|}{r})^3,} & \text{for } \frac{r}{2} < \| \mathbf{x} - \mathbf{p} \| \leq r, \\[8pt]
1, & \text{for } \| \mathbf{x} - \mathbf{p} \| > r.
\end{cases}
\end{equation}

\begin{wrapfigure}{l}[-3pt]{0.4\linewidth}
\vspace{-10pt}
\includegraphics[width=1.1\linewidth]{figures/Ilus_kernel_revision.pdf}
%\caption{Illustration of a domain with a discontinuity \EG{use $\Gamma$ instead of $\partial\Omega$}}
\vspace{-20pt}
\end{wrapfigure}
This formulation is derived from the cubic spline kernel, modified as one minus the original function to smooth out the singularity at zero distance. \Revision{We visualize the kernel shape on the right. The function is smooth over the entire domain, zero at $\| \mathbf{x} - \mathbf{p} \| = 0$ to cancel the singularity, and constant beyond the margin length $r$.} The final representation of $H$ is therefore: 
\begin{equation}
\label{eq:winding_para_smooth}
\begin{aligned}
    H(\mathbf{x}, \alpha) &= \int_{m}^{\alpha-m} \frac{\Gamma'(t) \cdot (\Gamma(t) - \mathbf{x})^\perp}{|\Gamma(t) - \mathbf{x}|^2} \, dt \\
   & +  K(x, \Gamma(0),  \|\Gamma(0) - \Gamma(m)\|) \int_{0}^{m} \frac{\Gamma'(t) \cdot (\Gamma(t) - \mathbf{x})^\perp}{|\Gamma(t) - \mathbf{x}|^2} \, dt \\
   & +  K(x, \Gamma(\alpha),  \|\Gamma(\alpha) - \Gamma(\alpha-m)\|) \int_{\alpha-m}^{\alpha} \frac{\Gamma'(t) \cdot (\Gamma(t) - \mathbf{x})^\perp}{|\Gamma(t) - \mathbf{x}|^2} \, dt.
\end{aligned}
\end{equation}

\section{Additional Details on Data-free Training}
\label{sec::appendix:training_details}
To construct skin\-ning eigen\-modes \cite{benchekroun2023FastComplemDynamics},
let $\Tilde{w}_\theta : \Omega \times \mathbb{R} \rightarrow \mathbb{R}$ be a continuous volumetric real-valued field, with a corresponding $w: \Omega \rightarrow \mathbb{R} : w = \Tilde{w}_\theta \circ \mathcal{L}$ planar real-valued skinning weight field discontinuous over $\Gamma$. Consider the $k=9$ dimensional reduced space parameterized by the matrix $Z \in \mathbf{R}^{3 \times 3}$, or equivalently by the vector $\bm{z} \in \mathbb{R}^{9} : \bm{z} = \textrm{flatten}(Z)$. Following \citet{benchekroun2023FastComplemDynamics}, the displacement field is 
\begin{align}
    \bm{u}(\bm{x}) = w(\bm{x}) \, Z\, \begin{bmatrix} \bm{x} \\ 1\end{bmatrix} \ .
\end{align}
Observe that $\bm{u}$ is linear in $Z$, therefore, if it were desired, it could be represented in the form $\bm{u} = \bm{z}^T \bm{\Phi}$ described above. Following~\citet{Modi:2024:Simplicits}, we uniformly stochastically sample the 9-dimensional handle $Z$ to train $\Tilde{w}_\theta$ by minimizing the elastic energy 
\begin{align} \label{eq:training-data-free}
    \mathcal{L}_{\text{data-free}} = E_{\text{elas}} = \int_{\Omega}\Psi(\bm{u}(\bm{x})) \, \mathrm{d}\bm{x} 
\end{align}
subject to $\langle w_\theta, 1 \rangle = 1$, i.e., orthogonality to a constant field, a constraint enforced by Gram–Schmidt orthogonalization \cite{chang2024neuralrepresentationshapedependentlaplacian}. We choose a St. Venant-Kirchhoff (StVK) material for the elastic energy density $\Psi$~\cite{barbivc2005real}. The domain integral is estimated using uniform stochastic quadrature. 

This process yields a 9-dimensional basis with a nonuniform weight function $w = \Tilde{w}_\theta \circ \mathcal{L}$. We extend this to an 18-dimensional basis by also incorporating a $3 \times 3$ matrix handle with uniform ($w=1$) weight~\cite{benchekroun2023FastComplemDynamics}.


\bibliographystyle{ACM-Reference-Format}
\bibliography{main}

\end{document}


